% ------------------------------------------------

\documentclass[12pt, a4paper]{report}

% ------------------------------------------------

% XeLaTex 的設定

% --------------------------
% 加這個就可以設定字體
\usepackage{fontspec}

% 使用xeCJK，其他的還有CJK或是xCJK
\usepackage[BoldFont]{xeCJK}

% Set english font (Use package <ifplatform> to check)
\usepackage{ifplatform}
\ifwindows
    % In Windows
    \setmainfont{Times New Roman}
\else
    % In *nix
    \usepackage{times}
\fi 

% 設定中英文的字型
% 字型的設定可以使用系統內的字型，而不用像以前一樣另外安裝
\setCJKmainfont{標楷體}

% 中文自動換行
\XeTeXlinebreaklocale "zh"

% 文字的彈性間距
\XeTeXlinebreakskip = 0pt plus 1pt

% 設定段落之間的距離
\setlength{\parskip}{0.3cm}

% 設定行距
\linespread{1.5}\selectfont

% ------------------------------------------------

% LaTex package

% --------------------------

\usepackage{./ncku/style/macros}
\usepackage{./ncku/style/utdiss}
\usepackage{graphicx}
\usepackage{latexsym}
\usepackage{hyperref}
\usepackage{pdfpages}

\usepackage{url}
\usepackage{multirow}
\newcommand{\tabincell}[2]{\begin{tabular}{@{}#1@{}}#2\end{tabular}}

\usepackage[sort&compress, sectionbib, super]{natbib}

%\usepackage{caption}
%\usepackage{enumitem}
%\usepackage{subcaption}

% ------------------------------------------------

% 有關學校對論文要求的設定

% --------------------------

% 學校浮水印 Watermark
\AddToShipoutPicture{
\put(0,0){
\parbox[b][\paperheight]{\paperwidth}{
\vfill
\centering
\includegraphics[]{./ncku/watersymbol.jpg}
\vfill
}}}

% --------------------------


% ----------------------------------------------------------------------------
%                           Command for cover page
% ----------------------------------------------------------------------------

% default variables definitions
% 此處只是預設值，不需更改此處
% 請更改 names.tex 內容
\newcommand\cTitle{}
\newcommand\eTitle{}
\newcommand\myCname{}
\newcommand\myEname{}
\newcommand\advisorCnameA{}
\newcommand\advisorEnameA{}
\newcommand\advisorCnameB{}
\newcommand\advisorEnameB{}
\newcommand\advisorCnameC{}
\newcommand\advisorEnameC{}
\newcommand\univCname{}
\newcommand\univEname{}
\newcommand\deptCname{}
\newcommand\fulldeptEname{}
\newcommand\deptEname{}
\newcommand\collEname{}
\newcommand\degreeCname{}
\newcommand\degreeEname{}
\newcommand\cYear{}
\newcommand\cMonth{}
\newcommand\eYear{}
\newcommand\eMonth{}
\newcommand\ePlace{}


% --------------------------

% 論文有關資料
%
% this file is encoded in utf-8
% v2.02 (Sep. 12, 2012)
% 填入你的論文題目、姓名等資料
% 如果題目內有必須以數學模式表示的符號，請用 \mbox{} 包住數學模式，如下範例
% 如果中文名字是單名，與姓氏之間建議以全形空白填入，如下範例
% 英文名字中的稱謂，如 Prof. 以及 Dr.，其句點之後請以不斷行空白~代替一般空白，如下範例
% 如果你的指導教授沒有如預設的三位這麼多，則請把相對應的多餘教授的中文、英文名
%    的定義以空的大括號表示
%    如，\renewcommand\advisorCnameB{}
%          \renewcommand\advisorEnameB{}
%          \renewcommand\advisorCnameC{}
%          \renewcommand\advisorEnameC{}

% 論文題目 (中文)
\renewcommand\cTitle{%
無線感測網路下的位置管理策略設計與效能評估
}

% 論文題目 (英文)
\renewcommand\eTitle{%
Design and Performance Evaluation of Mobility Management in Wireless Sensor Networks
}

% 我的姓名 (中文)
\renewcommand\myCname{胡庭瑜}

% 我的姓名 (英文)
\renewcommand\myEname{Ting-Yu Hu}

% 指導教授A的姓名 (中文)
\renewcommand\advisorCnameA{陳朝鈞 教授}

% 指導教授A的姓名 (英文)
\renewcommand\advisorEnameA{Dr.~Chao-Chun Chen}

% 指導教授B的姓名 (中文)
\renewcommand\advisorCnameB{陳朝鈞 教授}
%\renewcommand\advisorCnameB{}

% 指導教授B的姓名 (英文)
\renewcommand\advisorEnameB{Dr.~Chao-Chun Chen}

% 指導教授C的姓名 (中文)
\renewcommand\advisorCnameC{陳朝鈞 教授}
%\renewcommand\advisorCnameC{}

% 指導教授C的姓名 (英文)
\renewcommand\advisorEnameC{Dr.~Chao-Chun Chen}

% 校名 (中文)
\renewcommand\univCname{國立成功大學}

% 校名 (英文)
\renewcommand\univEname{National Cheng Kung University}

% 系所名 (中文)
\renewcommand\deptCname{資訊工程學系}

% 系所全名 (英文)
\renewcommand\fulldeptEname{Department of Computer Science and Information Engineering}

% 系所短名 (英文, 用於書名頁學位名領域)
\renewcommand\deptEname{CSIE}

% 學院英文名 (如無，則以空的大括號表示)
\renewcommand\collEname{College of Electrical Engineering and Computer Science}

% 學位名 (中文)
\renewcommand\degreeCname{碩士}

% 學位名 (英文)
\renewcommand\degreeEname{Master of Engineering}

% 口試年份 (中文、民國)
\renewcommand\cYear{一○一}

% 口試月份 (中文)
\renewcommand\cMonth{七}

% 口試年份 (阿拉伯數字、西元)
\renewcommand\eYear{2012}

% 口試月份 (英文)
\renewcommand\eMonth{July}

% 學校所在地 (英文)
\renewcommand\ePlace{Tainan, Taiwan}

%畢業級別；用於書背列印；若無此需要可忽略
\newcommand\GraduationClass{100}

%%%%%%%%%%%%%%%%%%%%%%


% --------------------------

% 使用 hyperref 在 pdf 簡介欄裡填入相關資料
\ifx \hypersetup \undefined
    % do nothing
	\relax
\else
	\hypersetup
    {
        unicode     = true,
	    pdftitle    = {\eTitle (\cTitle)},
	    pdfauthor   = {\myEname (\myCname)},
        pdfcreator  = {\univEname},
%        pdfproducer = {\univEname},
        pdfsubject  = {},
        pdfkeywords = {},
    }
\fi

% ------------------------------------------------
